\chapter{Úvod}
\label{chapter:Introduction}

Od sprístupnenia služby ChatGPT koncom roka 2022 sa generatívne nástroje umelej inteligencie stali bežnou súčasťou práce žiakov. Táto zmena predbehla vzdelávaciu politiku aj výskum: školy museli na nástroje reagovať skôr, než vznikli dôkazy o tom, ako pôsobia. Na Slovensku vstupuje otázka do rozhodujúcej fázy~--- od marca 2026 je gramotnosť v oblasti umelej inteligencie súčasťou Štátneho vzdelávacieho programu pre základné školy, v školskom roku 2026/2027 sa má rozšíriť na stredné školy a plánované je sprístupnenie desiatok tisíc licencií pre pedagógov~\cite{minedu2025plan,drucker2026competence}. Otázka teda neznie, či sa nástroje do vyučovania dostanú, ale v akej podobe.

Odpoveď, ktorú by na túto otázku bolo možné oprieť o domáce dáta, neexistuje. Systematická rešerš 155 empirických štúdií~\cite{garzon2025systematic} nenašla zo Slovenska ani jednu prácu spĺňajúcu kritériá zaradenia. Zahraničné zistenia sú pritom rozporuplné natoľko, že z nich nemožno odvodiť praktické odporúčanie: v harvardskom experimente adaptívny tutor prekonal aktívne učenie~\cite{kestin2025aitutoring}, v tureckom experimente na stredných školách generatívny nástroj znížil výsledky na následnom teste o sedemnásť percent~\cite{bastani2025harm}, a priame porovnanie adaptívneho nástroja s generickým vyhľadávačom nenašlo v učebných ziskoch rozdiel~\cite{slijepcevic2025khanmigo}.

Táto práca tvrdí, že rozpor je zdanlivý. Vysvetľuje ho jedno rozlíšenie a jedna hypotéza. Rozlíšenie spočíva v tom, že „výsledky vzdelávania“ nie sú jedna veličina, ale tri: výkon počas hodiny, keď je nástroj po ruke; okamžité vybavenie bezprostredne po nej; a trvalé vybavenie s transferom po odstupe. Tieto tri sa môžu usporiadať v opačnom poradí. Hypotéza potom znie, že nástroj zvyšuje prvé dve takmer za každých okolností, pretože odstraňuje úsilie, kým jeho vplyv na tretiu závisí od toho, či interakcia žiaka s ním zostala generatívna~--- žiak formuloval vlastné riešenie či vysvetlenie~--- alebo sa stala offloadingová, keď delegoval samotný kognitívny úkon.

Hypotéza má jednu vlastnosť, ktorá ju robí testovateľnou a zároveň dosiaľ neoverenou: opiera sa o veličinu, ktorú doteraz nikto nemeral priamo. Gerlich~\cite{gerlich2025offloading} ju zisťoval sebahodnotiacimi škálami, teda nástrojom, ktorý je pri správaní podliehajúcom sociálnej žiaducnosti systematicky skreslený. Kosmyna a kolektív~\cite{kosmyna2025brain} ju merali elektroencefalografiou, ktorá je v školskej triede nepoužiteľná. Bastaniová a kolektív ju z dát neodvodzovali vôbec.

Cieľom práce je túto medzeru zaplniť. Prácu tvoria dva prepojené príspevky. Prvým je \emph{index epistemickej kvality promptov}~--- pozorovaná, automaticky vypočítateľná miera toho, nakoľko žiak pri práci s nástrojom skutočne premýšľal, odvodená z promptov, ktoré odoslal, a zachytených nástrojom vyvinutým v technickej časti práce. Druhým je návrh krížového experimentu na slovenských stredných školách, ktorý oddeľuje účinok pedagogického ohraničenia nástroja od účinku samotného modelu tým, že generickú a adaptívnu podmienku stavia nad \emph{ten istý model, tie isté úlohy a to isté rozhranie}.

Práca je členená nasledovne. Kapitola~\ref{chapter:sota} analyzuje súčasný stav so zameraním na metodológiu siedmich kľúčových empirických štúdií a končí formuláciou zjednocujúcej hypotézy a ôsmich preregistrovaných hypotéz. Kapitola~\ref{chapter:metodika} opisuje dizajn experimentu, merné nástroje, index PEQI, otázku identifikácie mechanizmu a štatistický plán vrátane analýzy sily testu odvodenej simuláciou. Kapitola~\ref{chapter:vysledky} spúšťa uzavretý analytický plán na dátach zo známeho generujúceho modelu a overuje, že plán dokáže obnoviť to, čo doň bolo vložené. Kapitola~\ref{chapter:problem} zasadzuje riešenie do inžinierskeho a spoločenského kontextu vrátane posúdenia rizík, reprodukovateľnosti a udržateľnosti. Kapitola~\ref{chapter:conclusion} zhŕňa prínosy, obmedzenia a ďalší postup.
